%% EXEMPLO FICTÍCIO — Atividade 3: projeto mecânico
\begin{atividade}{Invólucro isolado do estimulador, em impressão 3D}{
  tipo         = mecanico,
  modalidade   = PES,
  alternativa  = VPC,
  horas        = 35,
  inicio       = 2024-06,
  fim          = 2024-08,
  projeto      = {NEBULA, estimulação elétrica funcional para ciclismo},
  papel        = Responsável,
  supervisor   = {Rafael Fictício, supervisor do projeto NEBULA},
  registros    = {NBL-127; NRO-PRO-014-902; NRO-PRO-003-902},
  competencias = {I, III, IV},
  disciplinas  = {Desenho Técnico; Fenômenos de Transporte; Materiais Elétricos},
}

\begin{contexto}
A placa da revisão A viajava numa caixa plástica improvisada: abrir a tampa expunha a alta tensão,
não havia ventilação, e a caixa, presa ao quadro da bicicleta com abraçadeiras de nylon, soltou-se duas
vezes com a vibração da pedalada. Com a placa nova (Atividade~1), o invólucro passou a ser parte do
projeto, com requisitos próprios, registrados no design record: nenhuma parte de alta tensão acessível
sem ferramenta; temperatura interna de no máximo \SI{60}{\celsius} com \SI{35}{\celsius} de ambiente;
massa de até \SI{400}{\gram} com a placa e a bateria; fixação num tubo de \SI{35}{\milli\meter},
removível em menos de um minuto; acesso aos conectores e ao botão de parada; e fabricação na impressora
do LABBIO.
\end{contexto}

\begin{contribuicao}
Concebi o invólucro, comparei as alternativas, fiz o modelo e os desenhos, calculei a temperatura
interna e a abraçadeira, imprimi as peças, montei o conjunto e fiz os ensaios térmico e de vibração. A
fiação do botão de parada foi feita por Lucas Fictício.
\end{contribuicao}

\begin{desenvolvimento}
\fichatecnica{
  ferramenta    = {FreeCAD 0.21 (modelo e desenhos) e PrusaSlicer 2.7 (fatiamento)},
  materiais     = {PETG (base, tampa e abraçadeira); TPU 95A (calço); insertos de latão M3},
  fabricacao    = {Impressão 3D por extrusão (FDM) no LABBIO: camada de 0,2\,mm, 4 perímetros, 30\,\% de preenchimento},
  pecas         = {Base, tampa, abraçadeira e calço; 4 parafusos M3},
  carregamentos = {Massa de até 0,38\,kg sob vibração de até 5\,g; 2,1\,W dissipados dentro},
  desenhos      = {Anexos do design record \texttt{NRO-PRO-014-902}},
  arquivos      = {Drive da equipe, pasta NEBULA/mecânica/involucro-rev1},
}

\textbf{As alternativas.} Comparei três soluções numa matriz de decisão com pesos definidos com o
supervisor --- isolação (30\,\%), desempenho térmico (20\,\%), custo (20\,\%), encaixe na placa
(15\,\%) e prazo (15\,\%): uma caixa comercial de ABS usinada (nota 3,4 de 5), uma caixa de alumínio
usinada (3,1: boa para o calor, mas condutora e cara) e uma peça impressa em PETG (4,3). O PETG
ganhou pela isolação, pelo encaixe sob medida e pela temperatura de transição vítrea, perto de
\SI{80}{\celsius}, bem acima da do PLA.

\textbf{O modelo.} A base recebe a placa sobre quatro colunas com insertos de latão; a tampa tem
ranhuras em labirinto, que deixam o ar passar sem deixar um caminho reto até a alta tensão, e o furo
do botão de parada; a abraçadeira, com calço de TPU, prende o conjunto ao tubo
(Figura~\ref{fig:ex-involucro}). As paredes têm \SI{2,4}{\milli\meter}, múltiplas da largura de
extrusão, e as folgas de encaixe, \SI{0,3}{\milli\meter}.

\begin{figure}[htbp]
  \centering
  % Vista explodida do invólucro (exemplo fictício), em projeção oblíqua:
% a base com a abraçadeira, a placa e a tampa com as ranhuras.
\begin{tikzpicture}[x={(7.2mm,0mm)}, y={(3.1mm,2.2mm)}, z={(0mm,7.2mm)},
    line join=round, font=\scriptsize,
    aresta/.style={draw=nro-tinta, line width=0.45pt}]
  % uma caixa: canto (x,y,z), largura (x), profundidade (y), altura (z), cores
  \def\caixa#1#2#3#4#5#6#7{%
    \filldraw[aresta, fill=#7!70] (#1+#4,#2,#3) -- (#1+#4,#2+#5,#3) -- (#1+#4,#2+#5,#3+#6) -- (#1+#4,#2,#3+#6) -- cycle;
    \filldraw[aresta, fill=#7!45] (#1,#2,#3) -- (#1+#4,#2,#3) -- (#1+#4,#2,#3+#6) -- (#1,#2,#3+#6) -- cycle;
    \filldraw[aresta, fill=#7!20] (#1,#2,#3+#6) -- (#1+#4,#2,#3+#6) -- (#1+#4,#2+#5,#3+#6) -- (#1,#2+#5,#3+#6) -- cycle;}
  % a abraçadeira, sob a base
  \filldraw[aresta, fill=nro-cinza!30] (4.2,4,-1.3) -- (7.8,4,-1.3) -- (7.8,4,0) -- (4.2,4,0) -- cycle;
  \draw[aresta, fill=white] (6,4,-1.9) circle (1.9mm);
  \node[anchor=west, align=left] at (13.2,4,-1.2) {abraçadeira para tubo de 35\,mm,\\braço de 8\,mm e calço de TPU};
  % a base (12 x 9 x 1,8)
  \caixa{0}{0}{0}{12}{9}{1.8}{nro-fundo}
  \foreach \x/\y in {1/1,11/1,1/8,11/8} \fill[nro-synapse!80!black] (\x,\y,1.8) circle (0.9mm);
  \node[anchor=west, align=left] at (13.2,4.5,0.9) {base em PETG, paredes de 2,4\,mm,\\insertos de latão M3};
  % as linhas de montagem (antes da placa e da tampa, que as escondem)
  \foreach \x/\y in {1/1,11/1,1/8,11/8} \draw[nro-claro, densely dashed] (\x,\y,1.8) -- (\x,\y,7.8);
  % a placa (10 x 7,5 x 0,2), com os componentes altos
  \caixa{1}{0.75}{4.2}{10}{7.5}{0.25}{nro-axon}
  \caixa{2}{2}{4.45}{2.5}{2.5}{0.6}{nro-cinza}
  \caixa{6}{1.2}{4.45}{4}{1.2}{0.7}{nro-tinta}
  \caixa{9.6}{3}{4.45}{1.2}{3}{0.9}{nro-synapse}
  \node[anchor=west, align=left] at (13.2,4.5,4.6) {placa de potência, revisão B\\(Atividade 1)};
  % a tampa, com as ranhuras em labirinto
  \caixa{0}{0}{7.8}{12}{9}{1.6}{nro-fundo}
  \foreach \k in {0,...,5} \filldraw[draw=nro-tinta, line width=0.3pt, fill=nro-tinta!55]
    (2+\k*1.3,5,9.4) -- (2.6+\k*1.3,5,9.4) -- (2.6+\k*1.3,8,9.4) -- (2+\k*1.3,8,9.4) -- cycle;
  \draw[aresta, fill=nro-pulso!80] (3.2,2.2,9.4) circle (2.2mm);
  \node[anchor=west, align=left] at (13.2,4.5,8.7) {tampa em PETG com ranhuras em\\labirinto e o botão de parada};
\end{tikzpicture}

  \caption{Vista explodida do invólucro: a base com a abraçadeira, a placa e a tampa.}
  \fonte{Elaborado pela autora.}
  \label{fig:ex-involucro}
\end{figure}

\textbf{A temperatura.} A placa dissipa \SI{2,1}{\watt} dentro do invólucro. Em convecção natural,
com $h \approx \SI{6}{\watt\per\square\meter\per\kelvin}$ e a área externa de
$A = \SI{0,034}{\square\meter}$, a elevação de temperatura estimada é
\begin{equation}
  \Delta T = \frac{P}{h A} = \frac{\SI{2,1}{\watt}}{\SI{6}{\watt\per\square\meter\per\kelvin} \times \SI{0,034}{\square\meter}} \approx \SI{10}{\kelvin},
  \label{eq:ex-termico}
\end{equation}
o que dá \SI{45}{\celsius} com \SI{35}{\celsius} de ambiente, dentro do requisito.

\textbf{A abraçadeira.} Tomei a massa de \SI{0,38}{\kilo\gram} sob 5\,g de vibração
($F = m a \approx \SI{19}{\newton}$) e dimensionei para $F = \SI{50}{\newton}$. O braço da abraçadeira
é uma viga em balanço de $L = \SI{30}{\milli\meter}$ e seção $b \times h$, com a tensão máxima
\begin{equation}
  \sigma = \frac{M c}{I} = \frac{F L \cdot h/2}{b h^3/12} = \frac{6 F L}{b h^2}.
  \label{eq:ex-viga}
\end{equation}
Com $b = \SI{12}{\milli\meter}$ e $h = \SI{5}{\milli\meter}$, a primeira versão dava
$\sigma = \SI{30}{\mega\pascal}$ --- fator de segurança 1,0 contra os \SI{30}{\mega\pascal} que
estimei para o PETG impresso na direção entre camadas (60\,\% dos \SI{50}{\mega\pascal} do material).
Aumentei para $h = \SI{8}{\milli\meter}$: $\sigma = \SI{11,7}{\mega\pascal}$ e fator 2,6.
\end{desenvolvimento}

\begin{resultados}
O conjunto com a placa e a bateria pesou \SI{352}{\gram}. No ensaio térmico, em estimulação contínua
na amplitude máxima, a primeira impressão, sem ranhuras, passou de \SI{49}{\celsius} com
\SI{26}{\celsius} de ambiente; com as ranhuras em labirinto, a temperatura estabilizou em
\SI{38}{\celsius}, uma elevação de \SI{12,5}{\kelvin}, perto dos \SI{10}{\kelvin} estimados em
\eqref{eq:ex-termico} (Figura~\ref{fig:ex-termico}). Extrapolada para \SI{35}{\celsius} de ambiente,
fica em \SI{48}{\celsius}, abaixo dos \SI{60}{\celsius} do requisito.

\begin{figure}[htbp]
  \centering
  % Temperatura do ar dentro do invólucro em estimulação contínua na
% amplitude máxima, a 26 °C de ambiente (exemplo fictício).
\begin{tikzpicture}
  \begin{axis}[
      width=0.86\linewidth, height=58mm,
      xlabel={Tempo (min)}, ylabel={Temperatura interna (\si{\celsius})},
      xmin=0, xmax=60, ymin=20, ymax=65,
      grid=major, grid style={nro-fundo},
      tick label style={font=\scriptsize}, label style={font=\footnotesize},
      legend pos=north west, legend cell align=left,
      legend style={font=\scriptsize, draw=nro-linha, fill=white, fill opacity=0.9, text opacity=1}]
    \addplot[only marks, mark=square*, mark size=1.4pt, nro-pulso-texto]
      coordinates {(0,25.8) (5,33.2) (10,38.1) (15,41.9) (20,44.3) (25,45.6) (30,46.8) (35,48.3) (40,48.4) (45,48.8) (50,49.7) (55,49.5) (60,49.9)};
    \addlegendentry{sem ranhuras (1.ª impressão)}
    \addplot[only marks, mark=*, mark size=1.5pt, nro-axon]
      coordinates {(0,26.1) (5,29.9) (10,32.6) (15,34.6) (20,35.7) (25,36.7) (30,36.7) (35,37.4) (40,38.1) (45,38.1) (50,38.5) (55,37.9) (60,38.3)};
    \addlegendentry{com ranhuras em labirinto}
    \addplot[domain=0:60, samples=61, nro-axon, thin] {26+12.5*(1-exp(-x/14))};
    \addlegendentry{ajuste: $\Delta T_\infty = 12{,}5$\,K, $\tau = 14$\,min}
    \addplot[domain=0:60, nro-cinza, densely dashed] {60};
    \addlegendentry{limite de projeto (60\,°C)}
  \end{axis}
\end{tikzpicture}

  \caption{Temperatura do ar dentro do invólucro em estimulação contínua na amplitude máxima, a
    \SI{26}{\celsius} de ambiente.}
  \fonte{Relatório de execução \texttt{NRO-PRO-003-902}.}
  \label{fig:ex-termico}
\end{figure}

No ensaio de vibração --- três horas na bicicleta, na bancada de rolo ---, nada se soltou, e a
remoção do conjunto levou \SI{40}{\second}. Dois problemas apareceram: os cantos da primeira base
empenaram na impressão, resolvido com borda de adesão e a impressora fechada; e o encaixe da tampa
ficou apertado demais, resolvido ao passar a folga de \num{0,2} para \SI{0,3}{\milli\meter}.
\end{resultados}

\begin{evidencias}
\evidencia{Arquivo de projeto}{involucro-rev1 (FreeCAD e STL)}{O modelo, os desenhos e os arquivos de impressão}{Drive da equipe (acesso a pedido)}
\evidencia{Arquivo controlado}{NRO-PRO-014-902}{O design record do invólucro, com a matriz de decisão e os cálculos}{SOMA › Arquivos}
\evidencia{Arquivo controlado}{NRO-PRO-003-902}{O relatório dos ensaios térmico e de vibração}{SOMA › Arquivos}
\evidencia{Atividade}{NBL-127}{O cartão do invólucro, atribuído a mim e concluído}{SOMA › Atividades › NEBULA}
\end{evidencias}

\begin{aprendizados}
Uma estudante de Engenharia Elétrica projetando uma peça mecânica: usei o Desenho Técnico mais do
que esperava, e a conta de convecção de Fenômenos de Transporte decidiu as ranhuras. O que aprendi fora
da sala: a peça impressa não é o material do catálogo --- a anisotropia entre camadas mudou o
dimensionamento ---, e um requisito de segurança (nenhuma parte de alta tensão acessível) vira
geometria, não só texto. Da próxima vez, eu imprimiria um corpo de prova para medir a resistência entre
camadas, em vez de estimá-la.
\end{aprendizados}
\end{atividade}
